\section{September 14, 2026}

\subsection{Direct sums and injectivity}

Let $V_1,\ldots,V_m$ be subspaces of a vector space $V$. Define the
summation map
\[
  \Sigma\colon V_1\times\cdots\times V_m
  \longrightarrow V_1+\cdots+V_m
\]
by
\[
  \Sigma(v_1,\ldots,v_m)=v_1+\cdots+v_m.
\]
This map is linear and surjective.

\begin{proposition}
  The sum $V_1+\cdots+V_m$ is direct if and only if the summation map
  $\Sigma$ is injective.
\end{proposition}

\begin{proof}
  The map $\Sigma$ is injective precisely when
  \[
    \Sigma(v_1,\ldots,v_m)=\Sigma(u_1,\ldots,u_m)
  \]
  implies
  \[
    (v_1,\ldots,v_m)=(u_1,\ldots,u_m).
  \]
  Equivalently, every vector in $V_1+\cdots+V_m$ has only one
  representation as a sum of vectors from the respective subspaces.
  This is exactly the definition of a direct sum.
\end{proof}

\subsection{Quotient vector spaces}

Let $U$ be a subset of $V$. For $v\in V$, define
\[
  v+U=\{v+u:u\in U\}.
\]
The set $v+U$ is called a \emph{translate} of $U$. When $U$ is a
subspace, a translate is also called a \emph{coset} of $U$.

\begin{example}
  Let
  \[
    U=\left\{\begin{bmatrix}x\\2x\end{bmatrix}:x\in\bb{R}\right\}
    \subseteq\bb{R}^2.
  \]
  Then
  \[
    \begin{bmatrix}2\\1\end{bmatrix}+U
    =\left\{\begin{bmatrix}2+x\\1+2x\end{bmatrix}:x\in\bb{R}\right\}.
  \]
  The original subspace is the line $y=2x$, while this translate is the
  parallel line $y=2x-3$.
\end{example}

\begin{definition}[Quotient space]
  If $U$ is a subspace of $V$, the \emph{quotient space} of $V$ by $U$
  is the set of all cosets of $U$:
  \[
    V/U=\{v+U:v\in V\}.
  \]
\end{definition}

For the subspace $U\subseteq\bb{R}^2$ in the preceding example,
$\bb{R}^2/U$ is the collection of all lines in $\bb{R}^2$ with slope
$2$.

\begin{proposition}\label{prop:coset-equivalence}
  Let $U$ be a subspace of $V$ and let $v,w\in V$. The following are
  equivalent:
  \begin{enumerate}[label=(\roman*)]
    \item $v-w\in U$;
    \item $v+U=w+U$;
    \item $(v+U)\cap(w+U)\neq\varnothing$.
  \end{enumerate}
\end{proposition}

\begin{proof}
  Suppose first that $v-w\in U$. If $v+u\in v+U$, then
  \[
    v+u=w+(v-w+u)\in w+U,
  \]
  because $v-w+u\in U$. Thus $v+U\subseteq w+U$. Since
  $w-v=-(v-w)\in U$, the reverse inclusion follows in the same way, so
  $v+U=w+U$.

  If $v+U=w+U$, then the two cosets plainly have nonempty intersection.

  Finally, suppose $(v+U)\cap(w+U)\neq\varnothing$. There are
  $u_1,u_2\in U$ such that
  \[
    v+u_1=w+u_2.
  \]
  Therefore
  \[
    v-w=u_2-u_1\in U.
  \]
\end{proof}

Thus two cosets of a subspace are either equal or disjoint. We define
addition and scalar multiplication on $V/U$ by
\begin{align*}
  (v+U)+(w+U)&=(v+w)+U, \\
  \alpha(v+U)&=\alpha v+U.
\end{align*}
Proposition~\ref{prop:coset-equivalence} shows that these operations do
not depend on the choice of representatives. With these operations,
$V/U$ is a vector space over $\bb{F}$, whose zero vector is the coset
$U=0+U$.

\subsection{Finite-dimensional vector spaces}

\begin{definition}[Finite-dimensional vector space]
  A vector space $V$ is \emph{finite-dimensional} if some finite
  collection of vectors spans $V$. It is \emph{infinite-dimensional}
  if it is not finite-dimensional.
\end{definition}

\begin{example}
  The coordinate space $\bb{F}^n$ is finite-dimensional because
  \[
    \bb{F}^n=\operatorname{span}\{e_1,\ldots,e_n\}.
  \]

  Likewise, the space of polynomials of degree at most $m$ is
  finite-dimensional because
  \[
    \mathcal P_m(\bb{F})
    =\operatorname{span}\{1,t,t^2,\ldots,t^m\}.
  \]
\end{example}

\begin{example}
  The vector space $\mathcal P(\bb{F})$ of all polynomials with
  coefficients in $\bb{F}$ is infinite-dimensional. Indeed, suppose a
  finite collection of polynomials spanned $\mathcal P(\bb{F})$, and let
  $m$ be the largest degree among the nonzero polynomials in the
  collection. Every polynomial in their span would have degree at most
  $m$, so $t^{m+1}$ could not belong to the span, a contradiction.
\end{example}
